\scanpage{221}
\mspage{113}
\deleted{4} Retour sur la condition d'isomorphie.

% slanted
\margin{(on suppose $R$ séparant en $S^\vee$)}
Soit
\[
 \deleted{\unsure{G}}\quad A=\mathrm{Aut}(\Gamma),\qquad
 \Gamma=(S,S^\vee,R) \tag*{(93)}
\]
\[
\begin{aligned}
 &Z\subset A\quad\text{un \unsure{sous-gr. ou|sous-groupe} commutatif donné}\\
 &\text{d'où}\quad\deleted{Z}\ \Pi_Z\subset\mathrm{Bij}(S,S^\vee),\ \text{sur lequel}\\
 &Z\text{ opère librement (a priori }\Pi_Z\text{ peut être vide, mais :)}
\end{aligned}\tag*{(94)}
\]
\[
 \text{un sous-}Z\text{-torseur }\pi\text{ de }\Pi_A,\tag*{(95)}
\]
On pose
\[
 G=\mathrm{Centr}_A(Z). \tag*{(96)}
\]
\[
 \deleted{\unsure{G'|\mathcal G'}=\{g\in G\mid\uline{\text{centralisant }Z}.}
 \tag*{\deleted{\unsure{(97)|(98)}}}
\]

\uline{Lemme} Soit $f\in\deleted{\unsure{G}}G$, Conditions équivalentes

(i) $\exists\ \varphi_0\in\pi$ tel que $f(A_{\varphi_0})$ de \deleted{\unsure{A}} la forme $A_\varphi$
avec $\varphi\in\pi$

(i bis) Pour \uline{tout} $\varphi_0\in\pi$, on a $f(A_{\varphi_0})$ de la forme $A_\varphi$
pour $\varphi=f(\varphi_0)$ convenable -- i.e. $f(\pi)=\pi$
(transport de structure)

\deleted{(ii) $\exists f\in$}
\deleted{L'ens. de ces $f$ forment un sous}

Les $f$ satisfaisant ces conditions \deleted{sont} $\in G$,
et forment un sous-groupe \inserted{$\unsure{\mathbf G|G}$} de $G$, contenant
\deleted{les} groupes \deleted{\unsure{Z_2|\mathcal A}} $\mathbf G_0=\mathrm{Aut}(S,A_{\varphi\deleted{\unsure{0}}})$ \deleted{pour}
\deleted{\unsure{$\varphi\in$} toute $\varphi\in\pi$, qui \unsure{est aussi}}
\inserted{(indépendant du choix de $\varphi$) et $Z$.} C'est aussi le sous-groupe
du groupe des autom. de la structure de
multigraphe à homothéties $(S,S^\vee,R,Z,\pi)$
induisant l'identité sur $Z$. On a un
diagramme, où la ligne est exacte
\diagram{\begin{tikzcd}[ampersand replacement=\&,column sep=large]
{\text{(97)}} \& 1\arrow[r] \& \mathbf G_0\arrow[r] \& \mathbf G\arrow[r,"\delta"]\arrow[d] \& Z\\
\& \& \& Z\arrow[ur,swap,"{\lambda\mapsto\lambda^2}"] \&
\end{tikzcd}}
% archive sheet number: 220

